\ifdefined\gkver\else\def\gkver{student}\fi
\documentclass[\gkver]{gaokaozhenti}
\begin{document}

\chapter{2025年安徽}

\begin{enumerate}
\item
%% number: 1
%% paperName: 2025年安徽高考第 1 题 4 分
%% typeId: 1
%% type: 单选题
%% chapter: 原子和原子核
%% point: 人工核反应
%% method: 无
%% score: 4
%% degree: 950
%% duplicateId: 0
%% body:
2025 年 4 月，位于我国甘肃省武威市的钍基熔盐实验堆实现连续稳定运行，标志着人类在第四代核电技术上迈出关键一步。该技术利用钍核（\ce{^{232}_{90}Th}）俘获 $x$ 个中子（\ce{^{1}_{0}n}），共发生 $y$ 次 $\beta$ 衰变，转化为易裂变的铀核（\ce{^{233}_{92}U}），则\xzanswer{B}

\fourchoices[answer=B]
{$x = 1$，$y = 1$}
{$x = 1$，$y = 2$}
{$x = 2$，$y = 1$}
{$x = 2$，$y = 2$}

%% answer: B
%% memo:
\memoanswer{
根据题意可知，钍核每俘获 1 个中子质量数加 1，电荷数不变，每发生一次 $\beta$ 衰变，质量数不变，电荷数加 1，钍核变成铀核，质量数加 1，电荷数加 2，则俘获 1 个中子，发生 2 次 $\beta$ 衰变，即 $x = 1$，$y = 2$。

故选 B。
}

\item
%% number: 2
%% paperName: 2025年安徽高考第 2 题 4 分
%% typeId: 1
%% type: 单选题
%% chapter: 机械波
%% point: 机械波的产生
%% method: 无
%% score: 4
%% degree: 750
%% duplicateId: 0
%% body:
如图，某同学演示波动实验，将一根长而软的弹簧静置在光滑水平面上，弹簧上系有一个标记物，在左端沿弹簧轴线方向周期性地推、拉弹簧，形成疏密相间的机械波。下列表述正确的是\xzanswer{D}

\onepicture[w=5.60cm]{figs/02.png}

\fourchoices[answer=D]
{弹簧上形成的波是横波}
{推、拉弹簧的周期越小，波长越长}
{标记物振动的速度就是机械波传播的速度}
{标记物由静止开始振动的现象表明机械波能传递能量}

%% answer: D
%% memo:
\memoanswer{
A．弹簧上形成的波的振动方向与传播方向平行是纵波，故 A 错误；

B．同一介质中，波的传播速度相同，则波的传播速度不变，推、拉弹簧的周期越小，波的周期越小，由公式 $\lambda = vT$ 可知，波长越短，故 B 错误；

C．标记物振动的速度反应的是标志物在平衡位置附近往复运动的快慢，机械波的传播速度是波在介质中的传播速度，二者不是同一个速度，故 C 错误；

D．标记物由静止开始振动，说明它获得了能量，这是因为机械波使得能量传递给标记物，则标记物由静止开始振动的现象表明机械波能传递能量，故 D 正确。

故选 D。
}

\item
%% number: 3
%% paperName: 2025年安徽高考第 3 题 4 分
%% typeId: 1
%% type: 单选题
%% chapter: 气体性质
%% point: 热力学第一定律
%% method: 无
%% score: 4
%% degree: 750
%% duplicateId: 0
%% body:
在恒温容器内的水中，让一个导热良好的气球缓慢上升。若气球无漏气，球内气体（可视为理想气体）温度不变，则气球上升过程中，球内气体\xzanswer{A}

\fourchoices[answer=A]
{对外做功，内能不变}
{向外放热，内能减少}
{分子的平均动能变小}
{吸收的热量等于内能的增加量}

%% answer: A
%% memo:
\memoanswer{
根据题意可知，气球缓慢上升的过程中，气体温度不变，则气体的内能不变，分子的平均动能不变，气体的体积变大，气体对外做功，由热力学第一定律可知，由于气体的内能不变，则吸收的热量与气体对外做的功相等。

故选 A。
}

\item
%% number: 4
%% paperName: 2025年安徽高考第 4 题 4 分
%% typeId: 1
%% type: 单选题
%% chapter: 直线运动
%% point: 多段匀变速直线运动
%% method: 无
%% score: 4
%% degree: 650
%% duplicateId: 0
%% body:
汽车由静止开始沿直线从甲站开往乙站，先做加速度大小为 $a$ 的匀加速运动，位移大小为 $x$；接着在 $t$ 时间内做匀速运动；最后做加速度大小也为 $a$ 的匀减速运动，到达乙站时速度恰好为 0。已知甲、乙两站之间的距离为 $8x$，则\xzanswer{A}

\fourchoices[answer=A]
{$x = \frac{1}{18}at^{2}$}
{$x = \frac{1}{16}at^{2}$}
{$x = \frac{1}{8}at^{2}$}
{$x = \frac{1}{2}at^{2}$}

%% answer: A
%% memo:
\memoanswer{
由题意可知，设匀加速直线运动时间为 $t'$，匀速运动的速度为 $v$，

匀加速直线运动阶段，由位移公式 $x = \frac{v}{2}t'$

根据逆向思维，匀减速直线运动阶段的位移等于匀加速直线运动阶段的位移，

则匀速直线运动阶段有 $8x - x - x = vt$

联立解得 $t' = \frac{t}{3}$

再根据 $x = \frac{1}{2}at'^{2}$

解得 $x = \frac{1}{18}at^{2}$

BCD 错误，A 正确。

故选 A。
}

\item
%% number: 5
%% paperName: 2025年安徽高考第 5 题 4 分
%% typeId: 1
%% type: 单选题
%% chapter: 运动和力
%% point: 连接体
%% method: 无
%% score: 4
%% degree: 550
%% duplicateId: 0
%% body:
如图，装有轻质光滑定滑轮的长方体木箱静置在水平地面上，木箱上的物块甲通过不可伸长的水平轻绳绕过定滑轮与物块乙相连。乙拉着甲从静止开始运动，木箱始终保持静止。已知甲、乙质量均为 $1.0 \Ukg$，甲与木箱之间的动摩擦因数为 $0.5$，不计空气阻力，重力加速度 $g$ 取 $10 \Umsq$，则在乙下落的过程中\xzanswer{C}

\onepicture[w=5.00cm]{\includesvg{figs/05}}

\fourchoices[answer=C]
{甲对木箱的摩擦力方向向左}
{地面对木箱的支持力逐渐增大}
{甲运动的加速度大小为 $2.5 \Umsq$}
{乙受到绳子的拉力大小为 $5.0 \UN$}

%% answer: C
%% memo:
\memoanswer{
A．因为物块甲向右运动，木箱静止，根据相对运动，甲对木箱的摩擦力方向向右，A 错误；

B．设乙运动的加速度为 $a$，只有乙有竖直向下的恒定加速度，

对甲、乙和木箱，由整体法，竖直方向受力分析有 $F_{N} = M_{\text{总}}g - ma$

则地面对木箱的支持力大小不变，B 错误；

CD．设绳子的弹力大小为 $T$，对甲受力分析有 $T - \mu mg = ma$

对乙受力分析有 $mg - T = ma$

联立解得 $a = 2.5 \Umsq$，$T = 7.5 \UN$

C 正确，D 错误。

故选 C。
}

\item
%% number: 6
%% paperName: 2025年安徽高考第 6 题 4 分
%% typeId: 1
%% type: 单选题
%% chapter: 机械振动
%% point: 简谐振动
%% method: 无
%% score: 4
%% degree: 550
%% duplicateId: 0
%% body:
在竖直平面内，质点 M 绕定点 O 沿逆时针方向做匀速圆周运动，质点 N 沿竖直方向做直线运动，M、N 在运动过程中始终处于同一高度。$t = 0$ 时，M、N 与 O 点位于同一直线上，如图所示。此后在 M 运动一周的过程中，N 运动的速度 $v$ 随时间 $t$ 变化的图像可能是\xzanswer{D}

\onepicture[w=3.60cm]{\includesvg{figs/06a}}

\fourchoices[answer=D, ispicture=true, h=2.40cm]
{figs/06b.png}
{figs/06c.png}
{figs/06d.png}
{figs/06e.png}

%% answer: D
%% memo:
\memoanswer{
因为 M、N 在运动过程中始终处于同一高度，所以 N 的速度 $v_{N}$ 与 M 在竖直方向的分速度 $v_{My}$ 大小相等，

设 M 做匀速圆周运动的角速度为 $\omega$，半径为 $r$，其竖直方向分速度 $v_{My} = \omega r \cos \omega t$

即 $v_{N} = \omega r \cos \omega t$

则 D 正确，ABC 错误。

故选 D。
}

\item
%% number: 7
%% paperName: 2025年安徽高考第 7 题 4 分
%% typeId: 1
%% type: 单选题
%% chapter: 磁场
%% point: 带电粒子在磁场中的运动
%% method: 无
%% score: 4
%% degree: 450
%% duplicateId: 0
%% body:
如图，在竖直平面内的 $Oxy$ 直角坐标系中，$x$ 轴上方存在垂直纸面向里的匀强磁场，磁感应强度大小为 $B$。在第二象限内，垂直纸面且平行于 $x$ 轴放置足够长的探测薄板 MN，MN 到 $x$ 轴的距离为 $d$，上、下表面均能接收粒子。位于原点 $O$ 的粒子源，沿 $Oxy$ 平面向 $x$ 轴上方各个方向均匀发射相同的带正电粒子。已知粒子所带电荷量为 $q$、质量为 $m$、速度大小均为 $\frac{qBd}{m}$。不计粒子的重力、空气阻力及粒子间的相互作用，则\xzanswer{C}

\onepicture[w=5.00cm]{figs/07.png}

\fourchoices[answer=C]
{粒子在磁场中做圆周运动的半径为 $2d$}
{薄板的上表面接收到粒子的区域长度为 $\sqrt{3}d$}
{薄板的下表面接收到粒子的区域长度为 $d$}
{薄板接收到的粒子在磁场中运动的最短时间为 $\frac{\pi m}{6qB}$}

%% answer: C
%% memo:
\memoanswer{
A．根据洛伦兹力提供向心力有 $qvB = \frac{mv^{2}}{R}$，可得 $R = \frac{mv}{qB} = d$，故 A 错误；

\onepicture[h=4.20cm]{figs/07_an.png}

B．当粒子沿 $x$ 轴正方向射出时，上表面接收到的粒子离 $y$ 轴最近，如图轨迹 1，根据几何关系可知 $s_{\text{上min}} = d$；当粒子恰能通过 N 点到达薄板上方时，薄板上表面接收点距离 $y$ 轴最远，如图轨迹 2，根据几何关系可知 $s_{\text{上max}} = \sqrt{3}d$，故上表面接收到粒子的区域长度为 $s_{\text{上}} = \sqrt{3}d - d$，故 B 错误；

C．根据图像可知，粒子可以恰好打到下表面 N 点；当粒子沿 $y$ 轴正方向射出时，粒子下表面接收到的粒子离 $y$ 轴最远，如图轨迹 3，根据几何关系此时离 $y$ 轴距离为 $d$，故下表面接收到粒子的区域长度为 $d$，故 C 正确；

D．根据图像可知，粒子恰好打到下表面 N 点时转过的圆心角最小，用时最短，有 $t_{\text{min}} = \frac{60^{\circ}}{360^{\circ}} \cdot \frac{2\pi m}{qB} = \frac{\pi m}{3qB}$，故 D 错误。

故选 C。
}

\item
%% number: 8
%% paperName: 2025年安徽高考第 8 题 4 分
%% typeId: 1
%% type: 单选题
%% chapter: 交流电
%% point: 变压器
%% method: 无
%% score: 4
%% degree: 450
%% duplicateId: 0
%% body:
某理想变压器的实验电路如图所示，原、副线圈总匝数之比 $n_{1}:n_{2} = 1:3$，A 为理想交流电流表。初始时，输入端 a、b 间接入电压 $u = 12\sqrt{2}\sin(100\pi t) \UV$ 的正弦式交流电，变压器的滑动触头 P 位于副线圈的正中间，电阻箱 $R$ 的阻值调为 $6 \UO$。要使电流表的示数变为 $2.0 \UA$，下列操作正确的是\xzanswer{B}

\onepicture[w=4.00cm]{figs/08.png}

\fourchoices[answer=B]
{电阻箱 $R$ 的阻值调为 $18 \UO$}
{副线圈接入电路的匝数调为其总匝数的 $\frac{1}{3}$}
{输入端电压调为 $u = 12\sqrt{2}\sin(50\pi t) \UV$}
{输入端电压调为 $u = 6\sqrt{2}\sin(100\pi t) \UV$}

%% answer: B
%% memo:
\memoanswer{
A．输入电压峰值为 $U_{m} = 12\sqrt{2} \UV$，则输入电压有效值为 $U_{1} = \frac{U_{m}}{\sqrt{2}} = 12 \UV$，滑动触头在正中间，根据变压比可知，输出电压 $U_{2} = \frac{0.5n_{2}}{n_{1}}U_{1} = 18 \UV$，若将电阻箱阻值调为 $18 \UO$，则电流为 $1 \UA$，故 A 错误；

B．若将副线圈匝数调为总匝数的 $\frac{1}{3}$，根据变压比可知，输出电压 $U_{2}' = \frac{\frac{1}{3}n_{2}}{n_{1}}U_{1} = 12 \UV$，则副线圈电流变为 $I_{2}' = \frac{U_{2}'}{R} = 2 \UA$，故 B 正确；

C．输入端电压调为 $u = 12\sqrt{2}\sin(50\pi t) \UV$ 时，其有效值不变，不会导致电流的变化，仍然为 $I = \frac{U_{2}}{R} = 3.0 \UA$，故 C 错误；

D．将输入电压峰值减小一半，则输入电压有效值变为 $U_{1}'' = \frac{\frac{1}{2}U_{m}}{\sqrt{2}} = 6 \UV$，输出电压 $U_{2}'' = \frac{0.5n_{2}}{n_{1}}U_{1}'' = 9 \UV$，副线圈电流变为 $I_{2}'' = \frac{U_{2}''}{R} = 1.5 \UA$，故 D 错误。

故选 B。
}

\item
%% number: 9
%% paperName: 2025年安徽高考第 9 题 5 分
%% typeId: 2
%% type: 多选题
%% chapter: 曲线运动
%% point: 开普勒定律
%% method: 无
%% score: 5
%% degree: 450
%% duplicateId: 0
%% body:
2025 年 4 月，我国已成功构建国际首个基于 DRO（远距离逆行轨道）的地月空间三星星座，DRO 具有“低能进入、稳定停泊、机动转移”的特点。若卫星甲从 DRO 变轨进入环月椭圆轨道，该轨道的近月点和远月点距月球表面的高度分别为 $a$ 和 $b$，卫星的运行周期为 $T$；卫星乙从 DRO 变轨进入半径为 $r$ 的环月圆形轨道，周期也为 $T$。月球的质量为 $M$，半径为 $R$，引力常量为 $G$。假设只考虑月球对甲、乙的引力，则\xzanswer{BC}

\fourchoices[answer=BC]
{$r = \frac{a + b + R}{2}$}
{$r = \frac{a + b}{2} + R$}
{$M = \frac{4\pi^{2}r^{3}}{GT^{2}}$}
{$M = \frac{4\pi^{2}R^{3}}{GT^{2}}$}

%% answer: BC
%% memo:
\memoanswer{
AB．对于题述环月椭圆轨道和环月圆轨道，根据开普勒第三定律有

$\frac{\left(\frac{a + b + 2R}{2}\right)^{3}}{T^{2}} = \frac{r^{3}}{T^{2}}$

可得 $r = \frac{a + b}{2} + R$

故 A 错误，B 正确；

CD．对于环月圆轨道，根据万有引力提供向心力可得 $\frac{GMm}{r^{2}} = m\left(\frac{2\pi}{T}\right)^{2}r$

可得 $M = \frac{4\pi^{2}r^{3}}{GT^{2}}$

故 C 正确，D 错误。

故选 BC。
}

\item
%% number: 10
%% paperName: 2025年安徽高考第 10 题 6 分
%% typeId: 2
%% type: 多选题
%% chapter: 电场
%% point: 带电粒子的平衡
%% method: 无
%% score: 6
%% degree: 350
%% duplicateId: 0
%% body:
如图，两个倾角相等、底端相连的光滑绝缘轨道被固定在竖直平面内，空间存在平行于该竖直平面水平向右的匀强电场。带正电的甲、乙小球（均可视为质点）在轨道上同一高度保持静止，间距为 $L$，甲、乙所带电荷量分别为 $q$、$2q$，质量分别为 $m$、$2m$，静电力常量为 $k$，重力加速度大小为 $g$。甲、乙所受静电力的合力大小分别为 $F_{1}$、$F_{2}$，匀强电场的电场强度大小为 $E$，不计空气阻力，则\xzanswer{ABD}

\onepicture[w=6.00cm]{\includesvg{figs/10}}

\fourchoices[answer=ABD]
{$F_{1} = \frac{1}{2}F_{2}$}
{$E = \frac{kq}{2L^{2}}$}
{若将甲、乙互换位置，二者仍能保持静止}
{若撤去甲，乙下滑至底端时的速度大小 $v = \sqrt{\frac{kq^{2}}{mL}}$}

%% answer: ABD
%% memo:
\memoanswer{
AB．对两球进行受力分析，设两球间的库仑力大小为 $F$，倾角为 $\theta$。

\onepicture[h=4.00cm]{figs/10_an.png}

对甲球根据平衡条件有 $F_{N1}\cos\theta = mg$，$F = F_{N1}\sin\theta + Eq$ ①

对乙球有 $F_{N2}\cos\theta = 2mg$，$F_{N2}\sin\theta = F + 2Eq$

联立解得 $F = 4Eq$ ②

故 $\frac{F_{1}}{F_{2}} = \frac{4Eq - Eq}{4Eq + 2Eq} = \frac{1}{2}$

同时有 $F = \frac{kq \cdot 2q}{L^{2}}$

解得 $E = \frac{kq}{2L^{2}}$

故 AB 正确；

C．若将甲、乙互换位置，若二者仍能保持静止，同理可得对甲有 $F_{N1}'\cos\theta = mg$，$F_{N1}'\sin\theta = F + Eq$

对乙有 $F_{N2}'\cos\theta = 2mg$，$F_{N2}'\sin\theta + 2Eq = F$

联立可得 $F + 4Eq = 0$，无解

假设不成立，故 C 错误；

D．若撤去甲，对乙球根据动能定理 $2mg \cdot \frac{L}{2}\tan\theta - 2Eq \cdot \frac{L}{2} = \frac{1}{2} \cdot 2mv^{2}$

根据前面分析由①②可知 $\tan\theta = \frac{3Eq}{mg}$

联立解得 $v = \sqrt{\frac{kq^{2}}{mL}}$

故 D 正确。

故选 ABD。
}

\item
%% number: 11
%% paperName: 2025年安徽高考第 11 题 6 分
%% typeId: 4
%% type: 实验题
%% chapter: 运动和力
%% point: 验证牛顿第二定律
%% method: 无
%% score: 6
%% degree: 550
%% duplicateId: 0
%% body:
某实验小组通过实验探究加速度与力、质量的关系。
\begin{enumerate}
\item
利用图~\ref{2025安徽11a} 装置进行实验，要平衡小车受到的阻力。平衡阻力的方法是：调整轨道的倾斜度，使小车 \tkanswer{c}。（选填正确答案标号）

a．能在轨道上保持静止

b．受牵引时，能拖动纸带沿轨道做匀速运动

c．不受牵引时，能拖动纸带沿轨道做匀速运动
\item
利用图~\ref{2025安徽11b} 装置进行实验，箱体的水平底板上安装有力传感器和加速度传感器，将物体置于力传感器上，箱体沿竖直方向运动。利用传感器测得物体受到的支持力 $F_{N}$ 和物体的加速度 $a$，并将数据实时传送到计算机。

①图~\ref{2025安徽11c} 是根据某次实验采集的数据生成的 $F_{N}$ 和 $a$ 随时间 $t$ 变化的散点图，以竖直向上为正方向。$t = 4 \Us$ 时，物体处于 \tkanswer{失重}（选填“超重”或“失重”）状态；以 $F_{N}$ 为横轴、$a$ 为纵轴，根据实验数据拟合得到的 $a-F_{N}$ 图像为图~\ref{2025安徽11d} 中的图线 a。

②若将物体质量增大一倍，重新进行实验，其 $a-F_{N}$ 图像为图~\ref{2025安徽11d} 中的图线 \tkanswer{d}。（选填“b”“c”或“d”）
\end{enumerate}

\fourpicture[labelA=2025安徽11a, labelB=2025安徽11b, labelC=2025安徽11c, labelD=2025安徽11d, numA=甲, numB=乙, numC=丙, numD=丁, wA=5.00cm, wB=3.00cm, wC=3.60cm, wD=3.40cm, align=right]
{figs/11a.png}
{figs/11b.png}
{figs/11c.png}
{figs/11d.png}

%% answer:
\jdanswer{
\begin{enumerate}
\item
c

\item
①失重　②d
\end{enumerate}
}
%% memo:
\memoanswer{
（1）平衡阻力的方法是：调整轨道的倾斜度，使小车不受牵引时，能拖动纸带沿轨道做匀速运动。

故选 c。

（2）①根据图像可知 $t = 4 \Us$ 时，加速度方向竖直向下，故处于失重状态；

②对物体根据牛顿第二定律 $F_{N} - mg = ma$

整理得 $a = \frac{1}{m} \cdot F_{N} - g$

可知图像的斜率为 $\frac{1}{m}$，故将物体质量增大一倍，图像斜率变小，纵轴截距不变，其 $a-F_{N}$ 图像为图丁中的图线 d。
}

\item
%% number: 12
%% paperName: 2025年安徽高考第 12 题 10 分
%% typeId: 4
%% type: 实验题
%% chapter: 电路
%% point: 多用表的使用
%% method: 无
%% score: 10
%% degree: 450
%% duplicateId: 0
%% body:
某同学设计了一个具有两种挡位（“$\times 1$”挡和“$\times 10$”挡）的欧姆表，其内部电路如图~\ref{2025安徽12a} 所示。电源为电池组（电动势 $E$ 的标称值为 $3.0 \UV$，内阻 $r$ 未知），电流表 G（表头）的满偏电流 $I_{g} = 20 \UmA$，内阻 $R_{g} = 45 \UO$，定值电阻 $R_{0} = 5 \UO$，滑动变阻器 $R$ 的最大阻值为 $200 \UO$。设计后表盘如图~\ref{2025安徽12b} 所示，中间刻度值为“15”。
\begin{enumerate}
\item
测量前，要进行欧姆调零：将滑动变阻器的阻值调至最大，闭合开关 $S_{1}$、$S_{2}$，此时欧姆表处于“$\times 1$”挡，将红表笔与黑表笔 \tkanswer{短接}，调节滑动变阻器的阻值，使指针指向 \tkanswer{0}（选填“0”或“$\infty$”）刻度位置。
\item
用该欧姆表对阻值为 $150 \UO$ 的标准电阻进行试测，为减小测量误差，应选用欧姆表的 \tkanswer{$\times 10$}（选填“$\times 1$”或“$\times 10$”）挡。进行欧姆调零后，将电阻接在两表笔间，指针指向图~\ref{2025安徽12b} 中的虚线位置，则该电阻的测量值为 \tkanswer{160} $\UO$。
\item
该同学猜想造成上述误差的原因是电源电动势的实际值与标称值不一致。为了测出电源电动势，该同学先将电阻箱以最大阻值（$9999 \UO$）接在两表笔间，接着闭合 $S_{1}$、断开 $S_{2}$，将滑动变阻器的阻值调到零，再调节电阻箱的阻值。当电阻箱的阻值调为 $228 \UO$ 时，指针指向“15”刻度位置（即电路中的电流为 $10 \UmA$）；当电阻箱的阻值调为 $88 \UO$ 时，指针指向“0”刻度位置（即电路中的电流为 $20 \UmA$）。由测量数据计算出电源电动势为 \tkanswer{2.8} $\UV$。（结果保留 2 位有效数字）
\end{enumerate}

\twopicture[labelA=2025安徽12a, labelB=2025安徽12b, numA=甲, numB=乙, wA=4.50cm, wB=6.00cm, align=right]{figs/12a.png}{figs/12b.png}

%% answer:
\jdanswer{
\begin{enumerate}
\item
短接，0

\item
$\times 10$，160

\item
2.8
\end{enumerate}
}
%% memo:
\memoanswer{
（1）测量前，要进行欧姆调零：将滑动变阻器的阻值调至最大，闭合开关 $S_{1}$、$S_{2}$，此时欧姆表处于“$\times 1$”挡，将红表笔与黑表笔短接，调节滑动变阻器的阻值，使指针指向 0 刻度位置。

（2）用该欧姆表对阻值为 $150 \UO$ 的标准电阻进行试测，为减小测量误差，应选用欧姆表的 $\times 10$ 挡；

进行欧姆调零后，将电阻接在两表笔间，指针指向图乙中的虚线位置，则该电阻的测量值为 $16 \times 10 \UO = 160 \UO$

（3）根据闭合电路欧姆定律 $E = I_{1}(r + R_{g} + R_{\text{阻}})$，$E = I_{2}(r + R_{g} + R_{\text{阻}}')$

代入数据 $I_{1} = 10 \UmA$，$I_{2} = 20 \UmA$

联立可得 $r = 7 \UO$，$E = 2.8 \UV$
}

\item
%% number: 13
%% paperName: 2025年安徽高考第 13 题 10 分
%% typeId: 6
%% type: 计算题
%% chapter: 光学
%% point: 全反射
%% method: 无
%% score: 10
%% degree: 450
%% duplicateId: 0
%% body:
如图，玻璃砖的横截面是半径为 $R$ 的半圆，圆心为 $O$ 点，直径与 $x$ 轴重合。一束平行于 $x$ 轴的激光，从横截面上的 $P$ 点由空气射入玻璃砖，从 $Q$ 点射出。已知 $P$ 点到 $x$ 轴的距离为 $\frac{\sqrt{2}}{2}R$，$P$、$Q$ 间的距离为 $\sqrt{3}R$。
\begin{enumerate}
\item
求玻璃砖的折射率；
\item
在该横截面沿圆弧任意改变入射点的位置和入射方向，使激光能在圆心 $O$ 点发生全反射，求入射光线与 $x$ 轴之间夹角的范围。
\end{enumerate}

\onepicture[w=5.50cm, align=right]{\includesvg{figs/13}}

%% answer:
\jdanswer{
\begin{enumerate}
\item
$\sqrt{2}$

\item
$(0, 45^{\circ}]$ 或 $[135^{\circ}, 180^{\circ})$
\end{enumerate}
}
%% memo:
\memoanswer{
（1）根据题意得出光路图如图所示。

\onepicture[h=4.00cm]{figs/13_an1.png}

根据几何关系可得 $\sin\alpha = \frac{\frac{\sqrt{2}}{2}R}{R}$，$\cos\gamma = \frac{\frac{\sqrt{3}}{2}R}{R}$，$\alpha = \beta$

可得 $\beta = 45^{\circ}$，$\gamma = 30^{\circ}$

根据折射定律 $n = \frac{\sin\alpha}{\sin\gamma} = \sqrt{2}$

（2）发生全反射的临界角满足 $\sin C = \frac{1}{n}$

可得 $C = 45^{\circ}$

\onepicture[h=3.80cm]{figs/13_an2.png}

要使激光能在圆心 $O$ 点发生全反射，激光必须指向 $O$ 点射入，如图所示。只要入射角大于 $45^{\circ}$，即可发生全反射，则使激光能在圆心 $O$ 点发生全反射，入射光线与 $x$ 轴之间夹角的范围 $(0, 45^{\circ}]$。由对称性可知，入射光线与 $x$ 轴之间夹角的范围还可以为 $[135^{\circ}, 180^{\circ})$。
}

\item
%% number: 14
%% paperName: 2025年安徽高考第 14 题 14 分
%% typeId: 6
%% type: 计算题
%% chapter: 曲线运动
%% point: 竖直平面圆周运动
%% method: 无
%% score: 14
%% degree: 450
%% duplicateId: 0
%% body:
如图，M、N 为固定在竖直平面内同一高度的两根细钉，间距 $L = 0.5 \Um$。一根长为 $3L$ 的轻绳一端系在 M 上，另一端竖直悬挂质量 $m = 0.1 \Ukg$ 的小球，小球与水平地面接触但无压力。$t = 0$ 时，小球以水平向右的初速度 $v_{0} = 10 \Ums$ 开始在竖直平面内做圆周运动。小球牵引着绳子绕过 N、M，运动到 M 正下方与 M 相距 $L$ 的位置时，绳子刚好被拉断，小球开始做平抛运动。小球可视为质点，绳子不可伸长，不计空气阻力，重力加速度 $g$ 取 $10 \Umsq$。
\begin{enumerate}
\item
求绳子被拉断时小球的速度大小，及绳子所受的最大拉力大小；
\item
求小球做平抛运动时抛出点到落地点的水平距离；
\item
若在 $t = 0$ 时，只改变小球的初速度大小，使小球能通过 N 的正上方且绳子不松弛，求初速度的最小值。
\end{enumerate}

\onepicture[w=5.50cm, align=right]{\includesvg{figs/14}}

%% answer:
\jdanswer{
\begin{enumerate}
\item
$v = 4\sqrt{5} \Ums$，$T = 17 \UN$

\item
$4 \Um$

\item
$2\sqrt{15} \Ums$
\end{enumerate}
}
%% memo:
\memoanswer{
（1）小球从最下端以速度 $v_{0}$ 抛出到运动到 M 正下方距离为 $L$ 的位置时，根据机械能守恒定律 $\frac{1}{2}mv_{0}^{2} = mg \cdot 2L + \frac{1}{2}mv^{2}$

在该位置时根据牛顿第二定律 $T - mg = m\frac{v^{2}}{L}$

解得 $v = 4\sqrt{5} \Ums$，$T = 17 \UN$

（2）小球做平抛运动时 $x = vt$，$2L = \frac{1}{2}gt^{2}$

解得 $x = 4 \Um$

（3）若小球经过 N 点正上方绳子恰不松弛，则满足 $mg = m\frac{v'^{2}}{2L}$

从最低点到该位置由动能定理 $\frac{1}{2}mv_{0}'^{2} = mg \cdot 5L + \frac{1}{2}mv'^{2}$

解得 $v_{0}' = 2\sqrt{15} \Ums$
}

\item
%% number: 15
%% paperName: 2025年安徽高考第 15 题 18 分
%% typeId: 6
%% type: 计算题
%% chapter: 电磁感应
%% point: 电磁感应能量分析
%% method: 无
%% score: 18
%% degree: 350
%% duplicateId: 0
%% body:
如图，平行光滑金属导轨被固定在水平绝缘桌面上，导轨间距为 $L$，右端连接阻值为 $R$ 的定值电阻。水平导轨上足够长的矩形区域 MNPQ 存在竖直向上的匀强磁场，磁感应强度大小为 $B$。某装置从 MQ 左侧沿导轨水平向右发射第 1 根导体棒，导体棒以初速度 $v_{0}$ 进入磁场，速度减为 0 时被锁定；从原位置再发射第 2 根相同的导体棒，导体棒仍以初速度 $v_{0}$ 进入磁场，速度减为 0 时被锁定，以此类推，直到发射第 $n$ 根相同的导体棒进入磁场。已知导体棒的质量为 $m$，电阻为 $R$，长度恰好等于导轨间距，与导轨接触良好（发射前导体棒与导轨不接触），不计空气阻力、导轨的电阻，忽略回路中的电流对原磁场的影响。求：
\begin{enumerate}
\item
第 1 根导体棒刚进入磁场时，所受安培力的功率；
\item
第 2 根导体棒从进入磁场到速度减为 0 的过程中，其横截面上通过的电荷量；
\item
从第 1 根导体棒进入磁场到第 $n$ 根导体棒速度减为 0 的过程中，导轨右端定值电阻 $R$ 上产生的总热量。
\end{enumerate}

\onepicture[w=7.00cm, align=right]{\includesvg{figs/15}}

%% answer:
\jdanswer{
\begin{enumerate}
\item
$\frac{B^{2}L^{2}v_{0}^{2}}{2R}$

\item
$\frac{mv_{0}}{BL}$

\item
$\frac{nmv_{0}^{2}}{2(n + 1)}$，$n = 1$，$2$，$3$，$\ldots$
\end{enumerate}
}
%% memo:
\memoanswer{
（1）第 1 根导体棒刚进入磁场时产生的感应电动势为 $E = BLv_{0}$

则此时回路的电流为 $I = \frac{E}{2R}$

此时导体棒受到的安培力 $F_{\text{安}} = BIL$

此时导体棒受安培力的功率 $P = F_{\text{安}}v_{0} = \frac{B^{2}L^{2}v_{0}^{2}}{2R}$

（2）第 2 根导体棒从进入磁场到速度减为 0 的过程中，根据动量定理有 $-B\bar{I}L \cdot \Delta t = 0 - mv_{0}$

其中 $\bar{I} \cdot \Delta t = q$

解得 $q = \frac{mv_{0}}{BL}$

（3）由于每根导体棒均以初速度 $v_{0}$ 进入磁场，速度减为 0 时被锁定，则根据能量守恒，每根导体棒进入磁场后产生的总热量均为 $Q = \frac{1}{2}mv_{0}^{2}$

第 1 根导体棒进入磁场到速度减为 0 的过程中，导轨右端定值电阻 $R$ 上产生的热量 $Q_{R1} = \frac{1}{2} \cdot Q$

第 2 根导体棒进入磁场到速度减为 0 的过程中，导轨右端定值电阻 $R$ 上产生的热量 $Q_{R2} = \frac{1}{2} \cdot \frac{1}{3} \cdot Q$

第 3 根导体棒进入磁场到速度减为 0 的过程中，导轨右端定值电阻 $R$ 上产生的热量 $Q_{R3} = \frac{1}{3} \cdot \frac{1}{4} \cdot Q$

第 $n$ 根导体棒进入磁场到速度减为 0 的过程中，导轨右端定值电阻 $R$ 上产生的热量 $Q_{Rn} = \frac{1}{n} \cdot \frac{1}{n + 1} \cdot Q$

则从第 1 根导体棒进入磁场到第 $n$ 根导体棒速度减为 0 的过程中，导轨右端定值电阻 $R$ 上产生的总热量 $Q_{R} = Q_{R1} + Q_{R2} + Q_{R3} + \ldots + Q_{Rn}$

通过分式分解和观察数列的“望远镜求和”性质，得出 $Q_{R} = \frac{n}{n + 1}Q = \frac{nmv_{0}^{2}}{2(n + 1)}$，$n = 1$，$2$，$3$，$\ldots$
}

\end{enumerate}

\end{document}
